%%%%%%%%%%%%%%%%%%%%%%%%%%%%%%%
%                             %
%      TEMPLATE "TESINA"      %
%       Andrea Salvador       %
%                             %
%%%%%%%%%%%%%%%%%%%%%%%%%%%%%%%


% CLASSE TESINA

% nota: la classe tesina include già i seguenti pacchetti:
%   babel, fontenc, graphicx, inputenc, setspace
% generalmente non servirà toglierli o modificarli

% opzioni per babel: italian, english
\documentclass[italian]{tesina-class/tesina}


% PACCHETTI GENERALI

% csquotes: virgolettati (\enquote{})
\usepackage[italian=quotes,english=american]{csquotes}
% link interni ed esterni
\usepackage{hyperref}
% gestione della bibliografia e dei riferimenti
\usepackage[style=numeric, backend=biber]{biblatex}
\addbibresource{bibliography.bib} % file dei riferimenti bibliografici


% PACCHETTI MATEMATICI

% ambienti equation, align etc
\usepackage{amsmath}
% notazione
\usepackage{amssymb}
% notazione
\usepackage{xfrac}


% istituzione (obbligatorio)
\renewcommand{\Istituzione}{Università degli Studi di Udine}
% dipartimento (facoltativo)
\renewcommand{\Dipartimento}{Scuola Superiore Universitaria \\ \enquote{Di Toppo-Wassermann}}
% altrimenti
%\renewcommand{\Dipartimento}{}
% corso di studi (facoltativo)
\renewcommand{\CorsoDiStudi}{}
% altrimenti
%\renewcommand{\CorsoDiStudi}{Laurea in \LaTeX}

% titolo
\renewcommand{\Titolo}{Titolo della tesina \\ anche \\ su \\ più \\ righe}

% in tutti i campi: se ci sono più persone separare con \and
% candidato (obbligatorio)
\renewcommand{\Candidato}{Candidato}
% relatore (obbligatorio)
\renewcommand{\Relatore}{Prof. Relatore}
% lasciare così se non ci sono correlatori
\renewcommand{\Correlatore}{}
% altrimenti
%\renewcommand{\Correlatore}{Dott. Correlatore 1 \and Dott. Correlatore 2}
% anno accademico (obbligatorio)
\renewcommand{\AnnoAccademico}{Anno Accademico 2026/2027}



\begin{document}

% frontespizio
\maketitle

% sommario
\begin{abstract}
Questo è un modello \LaTeX{} per una tesina della Scuola Superiore, semplice ma allo stesso tempo completo, riutilizzabile e adattabile.
È mirato a studenti di materie scientifiche, in particolar modo matematica e informatica, ma smanettando un po' si adatta a tutti.
\end{abstract}

%eventualmente in una nuova pagina
%\newpage

% indice
\tableofcontents



\newpage
% pagine numerate a partire da dopo l'indice
\pagenumbering{arabic}

\section{Introduzione}

Se questa repository o questo modello vi sono capitati tra le mani è probabile che vogliate scrivere una tesina\footnote{
    Altrimenti detta prova finale perché è l'esame finale di ogni anno nella Scuola Superiore Universitaria \enquote{Di Toppo-Wassermann} dell'Università degli Studi di Udine.
    Nulla vieta di usare questo modello per la \enquote{tesina di marzo} (prova intermedia), né per tesi, tesine o elaborati vari di altre istituzioni.
} e siate studenti di un corso scientifico, specialmente matematica e informatica.
Prevedendo questo uso i pacchetti predefiniti che accompagnano la classe \texttt{tesina} cercano di supportare i bisogni specifici di questi corsi di studi.
È sempre possibile aggiungere nuovi pacchetti \LaTeX{} ogni qual volta che si renda necessario.

Questo documento di esempio ha contemporaneamente lo scopo di fornire una classe \LaTeX{} per le tesine e di documentare la classe stessa.
Si assume un minimo di conoscenza del funzionamento elementare di \LaTeX{}, che se dovesse mancare si può sempre recuperare nei tutorial di Overleaf\footnote{
    A meno di variazioni dell'URL: \url{https://www.overleaf.com/learn}.
}

\subsection{Frontespizio}

Tutti gli elementi importanti del frontespizio sono modificabili agendo sui \verb|\renewcommand| del preambolo.
Con i comandi \verb|\Istituzione|, \verb|\Dipartimento| e \verb|\CorsoDiLaurea| (gli ultimi due facoltativi) si aggiustano le righe in alto.
Con \verb|\Titolo| si inserisce il titolo; sia il titolo sia ciò che c'è scritto sopra possono essere spezzati su più righe.

I comandi \verb|\Candidato|, \verb|\Relatore| e \verb|Correlatore| sono altrettanto autoesplicativi, notando che si possono inserire più nomi per ognuno di essi separandoli con il comando \verb|\and| per andare a capo (materialmente può avere senso solo nel caso dei correlatori).
Se il comando \verb|\Correlatore| è definito vuoto allora non viene mostrato nulla; gli altri due sono obbligatori.
Con il comando \verb|\AnnoAccademico| si può impostare ciò che viene scritto in fondo alla pagina.

L'immagine predefinita è il logo dell'Università di Udine.
Per cambiarlo bisogna agire sul file \verb|firstpage.cls| cambiando il percorso file con quello dell'immagine desiderata; nel caso non sappiate farlo, \emph{rangjaisi}.





\section{Pacchetti generali}

Descrizioni sintetiche dei pacchetti di uso comune inclusi di default.
Per una documentazione completa fate riferimento a quelle dei singoli pacchetti.

\subsection{\texttt{csquotes}}

\verb|csquotes| serve per formattare correttamente i virgolettati (comando \verb|\enquote|) in modo che \enquote{si chiudano bene}.

\subsection{\texttt{hyperref}}

\verb|hyperref| crea dei veri e propri link.
Vengono applicati in automatico a: rimandi nell'indice, rimandi alla bibliografia, note a piè di pagina. Si possono aggiungere link a siti web con \verb|\url|.

\subsection{\texttt{biblatex}}

\verb|biblatex| è fondamentale per la bibliografia.
Viene già incluso in questa classe con le seguenti opzioni: \texttt{style=numeric}, \texttt{backend=biber}.
Le citazioni appaiono così: \autocite{scrivere_una_tesina}.



\section{Pacchetti matematici}

In realtà anche questi sono di uso abbastanza ampio, utili in tutte le branche della scienza.

\subsection{\texttt{amsmath}}

\verb|amsmath| introduce ambienti per formule matematiche, in particolare \verb|equation|, \verb|align|, \verb|multiline|, \verb|split|.
C'è una gestione automatizzata delle equazioni numerate.
Oltre a questo, \verb|amsmath| introduce elementi di notazione matematica utili come \verb|\frac|, funzioni trigonometriche, sommatorie, integrali e operatori vari.
Sono importanti anche i font come \verb|\mathbb| ($\mathbb{R}$ etc.) e \verb|\mathcal| ($\mathcal{C}$ etc.).

\subsection{\texttt{amssymb}}

\verb|amssymb| aggiunge altra notazione matematica utile.

\subsection{\texttt{xfrac}}

\verb|xfrac| aggiunge le frazioni con notazione \enquote{inclinata} ($\sfrac{x}{y}$).



\newpage
\printbibliography

\end{document}